\documentclass[../main.tex]{subfiles}
\begin{document}

\resumeSubHeadingListStart

\resumeProjectHeading
{
  \titleItem{Better Pulse [WIP]} \quad [
    \href{https://github.com/bernardbdas/better-pulse}{\underline{Source
  Code}} ]   \\
  \emph{Backend: FastAPI, Pytorch, SQL, Ollama, Docker}
  \pipe
  \emph{Frontend: React Native, Expo}
}
{2026}
\resumeItemListStart
\resumeItem{Better Pulse (formerly Pulse-FL) is an attempt to bring
together edge intelligence and federated learning for healthcare domain.}
\resumeItem{Better Pulse will use the MIT-BIH ECG Datasets to train a
lightweight ECG model and finetune a pre-trained hugging face model.}
\resumeItem{The frontend is a React Native application that will be
  used to collect ECG data from users, train a local model and send the
updated weights to the backend for processing.}
\resumeItem{Will be implementing different federated aggregation
algorithms to find out which one is suitable for this usecase.}
\resumeItemListEnd

\resumeProjectHeading
{
  \titleItem{ECG Data Explorer} \quad [
    \href{https://github.com/bernardbdas/mit-bih}{\underline{Source
  Code}} ]   \\
  \emph{Tech Stack: wfdb, PyTorch, Streamlit, PEFT, Ollama, Python}
}
{2026}
\resumeItemListStart
\resumeItem{I used the MIT-BIH dataset to train a lightweight ECG
model and finetune pre-trained models like Random Forest and 1D ResNet.}
\resumeItem{The application is capable of analyzing the time-series
  data and bring out meaningful inferences using a local LLM like
gemma4 via Ollama.}
\resumeItem{The visualizations used have fine grained controls like
pause/playback and speed control using wfdb package.}
\resumeItem{The model is wrapped in an API that can be used to infer
  ECG data and return the predictions. This enables the user to bring
  their own ECG/EEG data in .csv file format to get better explanations
and inferences.}
\resumeItemListEnd

\resumeProjectHeading
{
  \titleItem{Mediapipe OpenCV Projects} \quad [
    \href{https://github.com/bernardbdas/practice-mediapipe}{\underline{Source
  Code}} ]   \\
  \emph{Tech Stack: Mediapipe, OpenCV, Streamlit, WebRTC, Python}
}
{2026}
\resumeItemListStart
\resumeItem{Developed a realtime object detection and landmarking
system using Mediapipe (newer Tasks API), OpenCV and Streamlit.}
\resumeItem{The application supports hand tracking, finger counting
and face detection.}
\resumeItem{The application can be accessed via a web interface using
WebRTC for video streaming.}
\resumeItem{The application now also supports pose estimation and
workout tracking .}
\resumeItemListEnd

\resumeProjectHeading
{
  \titleItem{Brain Image Segmentation} \quad [
    \href{https://github.com/bernardbdas/brain-tumor-detection.git}{\underline{Source
  Code}} ]   \\
  \emph{Tech Stack: Flask, TensorFlow, Keras, Numpy, Matplotlib}
}
{2026}
\resumeItemListStart
\resumeItem{Preprocessed MRI data with 3 tumor classes: 'meningioma',
'glioma', 'pituitary' and finetuned a pretrained VGG16 backbone.}
\resumeItem{The final model achieved 99.36\% accuracy on testing set,
after fine-tuning using .}
\resumeItem{The model can be accessed via a web interface using
Flask and can predict tumor type from MRI image.}
\resumeItem{WIP: Hosting the model using a NextJS frontend web
interface evolving into a fullstack brain image segementation app.}
\resumeItemListEnd

%  \resumeProjectHeading
%     {
%         \titleItem{Distributed Real-Time RAG \& Recommendation Engine}\pipe
%         \href{https://mediapipe-opencv-projects-701262404781.asia-south1.run.app/}{\underline{Live
% Demo}} \pipe
%         \href{https://github.com/bernardbdas/mediapipe-opencv-projects}{\underline{Source
% Code}} \pipe  \\
%         \emph{Tech Stack: PyTorch, Kafka, Spark, Triton,
% LangChain, Kubernetes}
%     }
%     {2026}
% \resumeItemListStart
%     \resumeItem{Engineered a scalable streaming pipeline using
% \textbf{Apache Kafka} and \textbf{PySpark} to process real-time
% user interaction logs and document metadata.}
%     \resumeItem{Trained and optimized a Two-Tower neural network in
% \textbf{PyTorch} for personalized content retrieval, tracking
% experiment lifecycles via \textbf{MLflow}.}
%     \resumeItem{Deployed the recommendation model via
% \textbf{NVIDIA Triton} (ONNX format) for low-latency inference,
% while orchestrating a concurrent \textbf{Hugging Face} LLM using
% \textbf{vLLM} and \textbf{LangChain} for document-grounded Q\&A.}
%     \resumeItem{Containerized all microservices with
% \textbf{Docker} and deployed the cluster using \textbf{Kubernetes}
% on cloud infrastructure to ensure high availability and auto-scaling.}
% \resumeItemListEnd

\resumeSubHeadingListEnd

\end{document}
